\documentclass[11pt,a4paper]{article}
\usepackage[utf8]{inputenc}
\usepackage{amsmath}
\usepackage{amsfonts}
\usepackage{amssymb}
\usepackage{graphicx}
\usepackage{hyperref}
\usepackage{geometry}
\geometry{a4paper, margin=1in}

\title{Progress Report: Escaping the Origin Trap via Adaptive Shot Boosting and the Transition to Hardware-Level Noise Mitigation}
\author{Quantum Nonlinear Dynamics Research Group}
\date{\today}

\begin{document}

\maketitle

\begin{abstract}
We report the successful stabilization of the Lorenz strange attractor on simulated quantum discrete architectures. The "Origin Trap"—a stochastic absorbing state induced by finite sampling—was successfully mitigated using a dynamic Adaptive Shot Boosting protocol ("The Quantum Microscope"). Moving forward, we transition from brute-force algorithmic variance reduction to structured Zero-Noise Extrapolation (ZNE) in the presence of physical depolarizing noise channels.
\end{abstract}

\section{Macroscopic Stabilization: Solving the Topological Floor}
In prior iterations, probability starvation near the unstable equilibrium origin $(0,0,0)$ caused finite measurement protocols to permanently round dynamic states to perfect zeroes, nullifying the chaotic repeller. 

To overcome this without losing topological coherence, we introduced the Quantum Resolution Floor ($p_{min} = \frac{0.5}{N_{shots}}$) paired with physical sign inertia. To prevent micro-orbit stagnation, we coupled this theoretical limit with an Adaptive Shot Boosting technique. When probability distributions fall into critical starvation zones ($< 10^{-4}$), our algorithm dynamically invokes a localized ``Quantum Microscope,'' scaling measurements from $50,000$ to $100,000$ shots for critical epochs. 

This hybrid threshold-boosting approach provided sufficient resolution to capture fractional probabilities and seamlessly reintegrate them into the Euler integrator. The result is the complete macroscopic topological tracing of the Lorenz butterfly lobes using purely Z-basis measurement tracking.

\section{The Scalability Ceiling and the Path to ZNE}
While Adaptive Shot Boosting analytically proves that the Statevector dynamics \textit{can} survive finite quantum sampling, the method relies on sheer computational brute force. Executing 100,000 shots on physical hardware for every step of an ODE integration represents an impractical overhead constraint. Furthermore, real NISQ hardware is plagued by inherent quantum decoherence beyond pure shot noise.

\subsection{Phase 2: Quantum Error Mitigation (OEM)}
For deployment on real processors, we must rely fundamentally on structured Quantum Error Mitigation (QEM). Our next evaluative phase explores \textbf{Zero-Noise Extrapolation (ZNE)}. 

Instead of hiding noise beneath exorbitant shot counts, ZNE utilizes artificial noise amplification (e.g., Unitary Folding $UU^\dag U$) to establish a noise propagation trend. By mapping the expectation values at multiple noise factors ($\lambda = 1, 3$), we use Richardson Extrapolation to retroactively infer the zero-error ($\lambda = 0$) theoretical probability.

We will now benchmark ZNE efficiency versus our Shot Boosting baseline to evaluate real-world feasibility.

\section{Control and Experimental Mitigation Results}
To establish the bounds of algorithmic resilience, we conducted a bipartite control experiment isolating statistical variance from physical decoherence.

\subsection{Phase 1: ZNE Control on Pure Shot Variance}
We first executed the Unitary Folding ($\lambda=1, 3$) extrapolated over a pristine simulator devoid of physical quantum error configurations, yielding purely 20,000 statistical \textit{shots} of noise. By applying standard Linear Richardson Extrapolation ($P_{ideal} = 1.5 P_1 - 0.5 P_3$), the simulation rapidly collapsed back into the Origin Trap. 

This experimentally verifies a crucial theoretical assumption: \textbf{ZNE is mathematically hostile towards pure statistical variance.} Because $P_1$ and $P_3$ are theoretically identical without depth-based physical decoherence, the subtractive extrapolation merely amplified the algorithmic sampling jitter.

\subsection{Phase 2: Exponential Extrapolation over Physical Device Noise}
Activating a calibrated structural Depolarizing Error mapping ($p=0.0001$), we effectively simulated deep physical gate-time degradation across the Unitary.

Because physical depolarizing channels decay exponentially relative to circuit depth (in contrast to linear statistical drift), we refactored the mitigation formula to an Exponential ZNE model: 
$$ P_{ideal} = P_1 \sqrt{\frac{P_1}{P_3}} $$
Enacting appropriate clipping bounds to prevent unphysical sub-floor division ($P_{ideal} \leq 1.0, P \geq p_{min}$), the algorithm successfully inverted the depolarization gradient. We retrieved striking continuous topological fidelity—reestablishing both macroscopic lobes of the strange attractor—while operating efficiently in a mere fraction of the computational magnitude required for pure brute-force ``Quantum Microscope'' shot augmentation.

\end{document}
